\section{Universal Chow generation supplies the missing forward
cycles}\label{universal-chow-generation-supplies-the-missing-forward-cycles}

\textbf{Scope.} The spectral target in CC-XC14 exists for every smooth
projective variety over a finite field; the remaining general gate is
the existence of algebraic maps from that variety to its target. Here a
geometric hypothesis supplies those maps. The input is Vial's published
decomposition theorem for rational Chow motives, not an additional Tate
conjecture for the variety under consideration. The finite-field descent
below is at the level of actual rational Chow correspondences. The final
conversion to the fixed target and the terminating search use the
abelian-sector inputs already recorded in AQ.15 and CC-XC11--15.

All Chow groups in this section have rational coefficients. A
\textbf{universal domain} \(\Omega\) is an algebraically closed field of
infinite transcendence degree over its prime field. In particular,
\(\overline{\mathbf F}_p\) is not such a domain. A Chow motive of
abelian type is an object of the full tensor subcategory closed under
finite sums, direct summands and duals generated by abelian varieties;
allowing finite étale schemes as additional generators gives its
abelian--Artin version. These are statements about rational equivalence,
before passage to homological or numerical equivalence.

\subsection{Primary input and
conventions}\label{primary-input-and-conventions}

Vial (2017, Theorem 3.11, pp.~213--214) works with covariant Chow
motives and writes \(\mathbf 1(1)\) for the Lefschetz motive. If a
morphism \(N\to h(X)\), with \(N\) finite-dimensional in the sense of
Kimura, surjects on algebraically trivial \(\mathrm{CH}_i\) over a
universal domain for \(i<l\), his self-dual conclusion splits \(h(X)\)
into a summand built from \(N,N^\vee\) and twisted curves, and a twisted
summand of a smooth projective variety of dimension at most
\(\dim X-2l\). This theorem has no characteristic-zero assumption. Its
special case with \(l=\lfloor\dim X/2\rfloor\) is Theorem 5,
pp.~197--198. The algebraically closed base-extension descent used below
is Vial's Theorem 1, p.~196, proved in §1.1, p.~199.

We use the contravariant convention of the preceding modules. Thus a
positive Lefschetz twist in Vial's convention is a negative étale Tate
twist here. The assertion that a motive belongs to the category
generated by twisted curves or abelian varieties is unchanged by this
translation. We retain all twists in the argument and do not assert an
untwisted Chow presentation from Vial's theorem.

\subsection{CC-XC16 --- A geometric sufficient condition, with actual
finite-field
descent}\label{cc-xc16-a-geometric-sufficient-condition-with-actual-finite-field-descent}

Let \(k=\mathbf F_q\), let \(X/k\) be smooth projective and
geometrically connected of pure dimension \(d\ge2\), and let
\(\Omega\supset\bar k\) be a universal domain. Put
\(l=\lfloor d/2\rfloor\). Suppose there is a finite direct sum \(N\) of
Tate twists of motives of products of smooth projective curves over
\(\Omega\), and a morphism of rational Chow motives

\[
                 u:N\longrightarrow h(X_\Omega),
 \tag{CC-XC.G1}
\]

written in Vial's covariant convention, such that

\[
 u_*:\mathrm{CH}_i(N)_{\mathrm{alg}}
       \longrightarrow \mathrm{CH}_i(X_\Omega)_{\mathrm{alg}}
 \quad\text{is onto for every }0\le i<l.
 \tag{CC-XC.G2}
\]

Here the finite direct sum and its finitely many correspondences are
part of the hypothesis. They may be defined over \(\Omega\), rather than
over \(k\). No finite-dimensionality assumption on \(h(X)\) is imposed.

Then \(h(X)\) belongs to the rational Chow category generated by abelian
varieties, finite étale \(k\)-schemes and Tate twists. In particular,
\(X\) satisfies the equivalent conditions CC-XC06 and CC-XC12. The
always-terminating construction CC-XC14 followed by the one-way search
CC-XC15 terminates for \(X\), and the accepted rational correspondences
give its fixed-target isomorphism for every \(\ell\ne p\).

\textbf{Complete deduction from the stated primary input.}

Every curve motive is finite-dimensional, and this property is preserved
by the finite sums, products and twists in \(N\); these are the
curve-source facts used in Vial's theorem and recalled on p.~195. Apply
Theorem 3.11 with algebraic equivalence and \(l=\lfloor d/2\rfloor\).
For negative \(i\), the target \(\mathrm{CH}_i(X_\Omega)\) is zero, so
(CC-XC.G2) gives the entire range \(i<l\) required there. Chow duality
gives the self-dual hypothesis for \(h(X_\Omega)\). The remaining
variety has dimension at most

\[
                    d-2\lfloor d/2\rfloor\in\{0,1\}.
 \tag{CC-XC.G3}
\]

Its motive, as well as every other summand produced by the theorem, is
generated by curves and Tate motives. Thus \(h(X_\Omega)\) is of abelian
type. This step uses an actual Chow splitting furnished by the cited
theorem, not an equality of Frobenius polynomials.

Apply Vial's Theorem 1 to the motive \(h(X_{\bar k})\) and the extension
\(\Omega/\bar k\). It gives a Chow splitting over \(\bar k\). Expanding
finite sums, tensor products and duals of abelian motives, we may write
it in the form

\[
 h(X_{\bar k})\ \xrightarrow{\ a\ }\
 T_{\bar k}=\bigoplus_\nu h(A_\nu)(n_\nu)
 \ \xrightarrow{\ b\ }\ h(X_{\bar k}),\qquad ba=1,
 \tag{CC-XC.G4}
\]

with abelian varieties \(A_\nu/\bar k\) and integers \(n_\nu\). Products
of abelian varieties are abelian, their dual motives are Tate twists of
their motives, and Tate motives can be included as summands of powers of
an elliptic curve, so this form loses no summands.

There is a finite extension \(k'/k\) over which all of (CC-XC.G4) is
defined \textbf{as a Chow splitting}. Indeed, the finitely many
varieties, their structure maps, the finitely many cycles representing
\(a,b\), and the finitely many rational-equivalence relations witnessing
\(ba=\Delta_X\) use finitely many elements of \(\bar k\). Each such
element lies in a finite extension of \(k\); their compositum is finite.
Enlarging it if necessary also defines the group structures and the
smooth projective models. Consequently

\[
 h(X_{k'})\xrightarrow{a'} T'
       \xrightarrow{b'}h(X_{k'}),\qquad b'a'=1
 \quad\text{in rational Chow motives over }k'.
 \tag{CC-XC.G5}
\]

For completeness, the descent from (CC-XC.G5) to membership over \(k\)
is as follows. Put \(e=[k':k]\). Regard a smooth projective
\(k'\)-scheme as a smooth projective \(k\)-scheme. A correspondence on
\(Y\times_{k'} Z\) extends by zero along the open-and-closed immersion
\(Y\times_{k'} Z\hookrightarrow Y\times_k Z\). This operation preserves
rational equivalence, identities and composition. To check composition,
extend to \(\bar k\): the schemes split into components indexed by the
embeddings \(k'\hookrightarrow\bar k\), and the induced correspondence
is block diagonal, with the same conjugate correspondence in each
matching block. The composition is the same block-diagonal composition.
The verification descends to rational Chow groups because base change on
rational Chow groups is injective (Vial, Lemma 1.2, p.~200). Denote this
additive operation by \(\operatorname{Ind}_{k'/k}\).

For an abelian variety \(A/k'\), let \(B=\operatorname{Res}_{k'/k}A\).
This is an abelian variety over \(k\), and

\[
 B_{k'}\simeq\prod_{\sigma\in\operatorname{Gal}(k'/k)}A^\sigma .
 \tag{CC-XC.G6}
\]

These precise restriction-of-scalars properties are given in Milne
(1972, §1, pp.~177--178). Projection onto the identity factor and
insertion into that factor, using the zero sections in the other
factors, have composite \(1_A\). Their graph correspondences therefore
exhibit

\[
 \operatorname{Ind}_{k'/k}h(A)
 \quad\text{as a Chow direct summand of}\quad
 h(B)\otimes h(\operatorname{Spec}k').
 \tag{CC-XC.G7}
\]

The right side is explicitly an abelian motive tensored with an Artin
motive. The statement does not replace this product by a single abelian
variety. Twists commute with the induction operation. Thus
\(\operatorname{Ind}_{k'/k}T'\), and its direct summand
\(h(X_{k'}\text{ viewed over }k)\), lie in the stated Chow subcategory.

Finally, the finite étale projection \(\tau:X_{k'}\to X\), with source
now viewed over \(k\), supplies pullback and trace correspondences
satisfying

\[
                         \tau_*\tau^*=e\,1_{h(X)} .
 \tag{CC-XC.G8}
\]

This is the graph-correspondence composition identity: the projection
from the corresponding fibre product onto the diagonal has degree \(e\).
It also follows from proper pushforward and the projection formula.
Division by the positive integer \(e\) is allowed over \(\mathbf Q\).
Hence \(\tau^*\) and \(e^{-1}\tau_*\) split \(h(X)\) off the induced
object, proving the claimed Chow inclusion over the specified field
\(k\).

Pass these actual correspondences to any homological realization.
CC-XC06 gives certification; CC-XC12 identifies the resulting
homological object with the fixed rational spectral target. CC-XC14
constructs a concrete untwisted numerical presentation of that target,
and CC-XC15 proves that its forward-cycle search terminates exactly when
a forward isomorphism exists. This proves the algorithmic conclusion. It
does not assert a Chow isomorphism to the particular numerical projector
output by CC-XC14. \(\square\)

\subsection{CC-XC17 --- Universal zero-cycles, including rationally
chain connected
threefolds}\label{cc-xc17-universal-zero-cycles-including-rationally-chain-connected-threefolds}

Let \(X/k\) be smooth projective and geometrically connected with
\(2\le d\le3\). If

\[
              \deg:\mathrm{CH}_0(X_\Omega)_{\mathbf Q}
                         \xrightarrow{\ \sim\ }\mathbf Q ,
 \tag{CC-XC.G9}
\]

then \(h(X)\) belongs to the same rational Chow abelian--Artin category
and its fixed-target forward-cycle search terminates. In fact, Vial's
theorem in this case expresses the Chow motive as a direct summand of a
finite sum of twisted curve motives. In particular the conclusion
applies when \(X\) is geometrically rationally chain connected, in every
characteristic.

\textbf{Complete proof.} Algebraically trivial zero-cycles have degree
zero. Thus (CC-XC.G9) implies
\(\mathrm{CH}_0(X_\Omega)_{\mathrm{alg}}=0\). Apply CC-XC16 with \(N=0\)
and \(l=1\). More precisely, one can apply Vial's Theorem 3.11 directly
to \(0\to h(X)\) over \(k\), since its Chow-group hypothesis is tested
over \(\Omega\): the self-dual remainder has dimension at most
\(d-2\le1\), and its other terms are twisted curves. This proves the
asserted stronger curve-source form with rational coefficients; no
numerical projector-lifting conjecture enters.

For the rational-chain-connected assertion, over an algebraically closed
extension such as \(\Omega\), any two points are joined by a connected
finite chain of images of \(\mathbf P^1\). On \(\mathbf P^1\), the
difference of two points is principal. Proper pushforward to
\(X_\Omega\) therefore makes the endpoints of each member of the chain
rationally equivalent; adding along the chain makes its two endpoints
rationally equivalent. Since closed points of \(X_\Omega\) are
\(\Omega\)-points, every zero-cycle is rationally equivalent to its
degree times one chosen point. This proves (CC-XC.G9), even integrally
over \(\Omega\). The standard rational-chain-connected implication and
its transfer formulation are also recorded in the published
Kahn--Sujatha Proposition 3.1.4(2), p.~400; Proposition 3.1.1,
pp.~398--399, gives the equivalent universal rational zero-cycle
formulations.

This argument uses smoothness and projectivity of \(X\), and uses
rational coefficients for the resulting motivic splitting. It does not
infer integral descent of a diagonal decomposition, or universal
integral triviality over nonclosed extensions of \(k\). The restriction
\(d\le3\) is essential to the application: if \(d\ge4\), vanishing of
only \(\mathrm{CH}_0(X_\Omega)_{\mathrm{alg}}\) leaves a possible
residual motive of dimension \(d-2\), which need not be generated by
curves. \(\square\)

\subsection{\texorpdfstring{Why the universal-domain condition cannot be
replaced by
\(\bar k\)}{Why the universal-domain condition cannot be replaced by \textbackslash bar k}}\label{why-the-universal-domain-condition-cannot-be-replaced-by-bar-k}

There is a direct test case. Let \(C/k\) be a smooth projective
geometrically connected curve of positive genus. Its Jacobian has only
torsion \(\bar k\)-points: each point is defined over a finite field,
whose group of rational points is finite. Consequently
\(\mathrm{CH}_0(C_{\bar k})_{\mathbf Q}^0=J(C)(\bar k)\otimes\mathbf Q=0\).
Nevertheless \(H^1(C)\ne0\). Moreover, after enlarging to a universal
domain \(\Omega\), the generic point of \(J(C)\) over \(\bar k\),
embedded in \(\Omega\), is not torsion: an equation \([n]x=0\) would
force the generic point to lie in the finite kernel of \([n]\). Thus
\(J(C)(\Omega)\otimes\mathbf Q\ne0\), and the same degree-zero Chow
group over \(\Omega\) is nonzero. The canonical degree-zero
Chow/Jacobian identification for a pointed curve gives these equalities;
a point exists over \(\bar k\).

A fixed-field vanishing therefore does not verify the universal
vanishing required by the theorem. This example distinguishes the
hypotheses; it is not a counterexample to abelian-type membership, since
a curve is itself of abelian type. Vial's zero-cycle discussion in §3.1,
pp.~207--208, explicitly explains why a test over a finite field cannot
replace the field-extension condition. In CC-XC16--17 the universal
Chow-generation assumption is geometric input; the earlier integral
comparison of Albanese relations in CC-ALB01--02 does not supply it
automatically.

\subsection{Resulting change to the original cote 15
gate}\label{resulting-change-to-the-original-cote-15-gate}

The forward-cycle existence problem is discharged for the stated
universal-Chow-generation class, including every smooth projective
geometrically rationally chain connected surface or threefold. The
conclusion is actual Chow membership, hence stronger than the required
numerical or homological inclusion. Neither the rational-chain-connected
test nor the more general Chow-generation condition is asserted for an
arbitrary input variety. For such an input, the forward-cycle existence
gate remains; the one-sided search still has no negative-membership
stopping rule.

\subsection{Primary bibliography and exact
locators}\label{primary-bibliography-and-exact-locators}

\begin{itemize}
\tightlist
\item
  Vial, C. (2017). Remarks on motives of abelian type. \emph{Tohoku
  Mathematical Journal, Second Series, 69}(2), 195--220.
  https://doi.org/10.2748/tmj/1498269623 . Published source:
  https://www.jstage.jst.go.jp/article/tmj/69/2/69\_1498269623/\_pdf/-char/en
  . Exact locators: conventions and curve finite-dimensionality, p.~195;
  Theorem 1, p.~196, and its §1.1 proof, p.~199; Theorem 5,
  pp.~197--198; Lemma 1.2, p.~200; universal-domain zero-cycle
  discussion, §3.1, pp.~207--208; Theorem 3.11 and complete proof,
  pp.~213--214. The journal version was received 24 October 2014 and
  revised 21 July 2015. The website's later release date is not a new
  theorem date.
\item
  Kahn, B., \& Sujatha, R. (2016). Birational motives, I: Pure
  birational motives. \emph{Annals of K-Theory, 1}(4), 379--440.
  https://doi.org/10.2140/akt.2016.1.379 . Published source:
  https://msp.org/akt/2016/1-4/akt-v1-n4-p02-s.pdf . Exact locators:
  Proposition 3.1.1, pp.~398--399; Proposition 3.1.4(2), p.~400,
  including its proof. The older preprint's Corollary 3.1.2 is not the
  published locator for rational chain connectedness.
\item
  Milne, J. S. (1972). On the arithmetic of abelian varieties.
  \emph{Inventiones Mathematicae, 17}, 177--190.
  https://www.jmilne.org/math/articles/1972a.pdf . Exact locator: §1,
  pp.~177--178, for restriction of scalars of an abelian variety and its
  product decomposition after finite separable base change. This is the
  same primary input already inspected for CC-XC08.
\end{itemize}
